
\begin{table}[htbp]
   
   \caption{\label{tab:6yr_huc02fe_inorg_ravalli_2005_k2} Effect of cumulative upstream coal production on Inorganic Chemicals, SYR2 1998--2005}
   \bigskip
   
   \centering
   
   \begin{adjustbox}{width = \textwidth, center}
      
      %start:tab
      
      \begin{tabular}{lcccc}
         \toprule
         
          & \multicolumn{4}{c}{Mean conc.}\\
         
                                             & Arsenic  & Nitrate        & Barium   & Selenium \\   
         
                                             & (1)      & (2)            & (3)      & (4)\\  
         
         \midrule 
         Cumul. upstream coal prod. (10M ST) & 0.0001   & 0.0883$^{***}$ & 0.0136   & 0.0004$^{**}$\\   
                                             & (0.0001) & (0.0257)       & (0.0112) & (0.0002)\\   
         
          \\
         
         Observations                        & 366      & 472            & 353      & 350\\  
         
         \bottomrule
         
      \end{tabular}
      
      %end:tab
      
      
   \end{adjustbox}
      {\tiny\linespread{1}\selectfont \par \raggedright 
   \textit{Notes:} Within each chemical, the column shows mean measured concentration from the EPA 6-Year Review. Non-detect values replaced by MDL$/\sqrt{2}$ following Ravalli et al.~(2022). Explanatory variable is cumulative coal production since 1985 (in 10 million short tons) in watersheds within two flow steps upstream of the utility's intake. Sample: utilities with a coal mine within two flow steps upstream of their intake and no coal mine colocated with their intake. Standard errors clustered at the utility level. All specifications include utility and HUC02 $\times$ year fixed effects. *** p$<$0.01, ** p$<$0.05, * p$<$0.1.}
   
\end{table}


